\section{A finite algebra of rational outcomes}
\label{sec:rational}

The interaction object admits several analyses. To give the logical
direction a concrete target, consider finite perfect-information trees
with one winner at each leaf and no strategic terminal-return stage.
Let $\langle j\rangle$ name the leaf won by $j$, and write
$G\oplus_r H$ for a binary choice owned by $r$.
An owned quantifier $\Q_r(x:A).G(x)$ is the corresponding choice over a
finite nonempty action type. Continuation trees may have different shapes.
Every player values winning at $1$ and every other outcome at $0$.
Strategies retain decisions at every history, as in
\cref{sec:policies}.

Let $B(G)\subseteq\Roles$ be the winners attained by full pure
subgame-perfect profiles. This is a specified analysis of the game, rather
than its definition. Its recursive calculation uses ordinary backward
induction, retaining all ties.

\begin{proposition}[Exact recursive calculation]
\label{prop:rational}
Every such tree has a nonempty set $B(G)$. For a leaf,
$B(\langle j\rangle)=\{j\}$. At a choice with child values
$S_a=B(G_a)$, put $U=\bigcup_a S_a$. Then
\[
 B(\Q_r a.G_a)=
 (\{r\}\cap U)\ \cup\
 \begin{cases}
 U\setminus\{r\},&
       \forall a,\ S_a\setminus\{r\}\ne\varnothing,\\
 \varnothing,&\text{otherwise}.
 \end{cases}
\]
\end{proposition}
\begin{proof}
Induct on the tree. A full profile is subgame perfect precisely when its
restriction to every child is subgame perfect and the root choice maximizes
the owner's payoff among those prescribed child outcomes.
For sufficiency of this characterization, a unilateral deviation within
a selected child cannot improve on that child's equilibrium outcome;
a deviation changing the root cannot improve on the maximum of these
outcomes. Necessity follows from the subgame condition and root deviations.

If some child can return winner $r$, choose that equilibrium there and
arbitrary equilibria elsewhere. Choosing that child is optimal for $r$.
A losing outcome $j\ne r$ can be chosen at the root precisely when every
child admits a losing equilibrium outcome: select one in each off-path
child, and the desired $j$ in its child. The root owner is then indifferent.
If a child admits only winner $r$, a losing root choice cannot be optimal.
These constructions also establish nonemptiness. Only finitely many
independent child equilibria must be assembled.
\end{proof}

For binary choice this simplifies to an operation on nonempty subsets:
\[
 S\mathbin{\star_r}T=
 \begin{cases}
 \{r\},&S=\{r\}\ \text{or}\ T=\{r\},\\
 S\cup T,&\text{otherwise}.
 \end{cases}
\]
Thus $B(G\oplus_r H)=B(G)\star_r B(H)$.
The singleton $\{r\}$ is absorbing: a child in which $r$ must win removes
all losing choices at the parent. A child merely admitting an $r$-win
does not have the same effect, because its off-path equilibrium can
prescribe a losing outcome.

\begin{lemma}[Algebraic laws and generated values]
\label{lem:rational-algebra}
Each $\star_r$ is associative, commutative, and idempotent.
Role permutations transport the operations:
$\pi(S\star_r T)=\pi S\star_{\pi r}\pi T$.
For $K=2$, the values generated by one-winner leaves are just the two
singletons. For $K\ge3$, every nonempty subset of $\Roles$ is generated.
\end{lemma}
\begin{proof}
For associativity, if any of three arguments equals $\{r\}$, both
bracketings yield $\{r\}$. Otherwise each intermediate union is different
from $\{r\}$ and both yield the three-way union. Commutativity,
idempotence, and equivariance follow from the displayed definition.

For $K=2$, every operation on two singleton values returns a singleton.
For $K\ge3$, realize $\{j,k\}$ by letting a third role choose between
leaves $j,k$. Inductively, to realize a set $S$ of size at least three,
choose $s\in S$ and an owner $r\in S\setminus\{s\}$.
Combine already realized $S\setminus\{s\}$ with $\{s\}$ under $r$.
Neither operand is $\{r\}$, so their value is $S$.
\end{proof}

At $K=2$, identify $\{0\}$ with true and $\{1\}$ with false.
Then $\star_0$ is disjunction and $\star_1$ is conjunction.
Their finite folds give the usual existential and universal readings,
and swapping the roles complements the value.
For three roles, indifference gives additional values and invalidates
classical distributivity. In particular,
\[
 B((\langle0\rangle\oplus_0\langle1\rangle)
                  \oplus_1\langle2\rangle)=\{0,2\},
\]
whereas
\[
 B((\langle0\rangle\oplus_1\langle2\rangle)
           \oplus_0(\langle1\rangle\oplus_1\langle2\rangle))
       =\{0,1,2\}.
\]
The repeated branch changes which off-path equilibrium can support an
indifferent choice. The calculation concerns full subgame-perfect
profiles, not local greediness along the executed path.

\begin{corollary}[No Boolean meet/join product representation]
\label{cor:no-boolean-product}
For $K\ge3$, this generated algebra has no injective homomorphism into a
product of Boolean algebras in which each owned operation acts in each
coordinate as either meet or join.
\end{corollary}
\begin{proof}
Any pair of such coordinate operations satisfies
$(a\star_r b)\star_s c
 =(a\star_s c)\star_r(b\star_s c)$, whether the coordinate operations
coincide or differ. Products and subalgebras preserve this identity.
The displayed three-role calculation refutes it in the generated algebra.
An injective homomorphism would preserve and reflect the equality.
\end{proof}

Consequently this rational-play analysis is not a faithful collection of
classical meet/join viewpoints. It offers a finite equational target:
terms have role-indexed choices, role permutations, and valuations in the
generated algebra; validity of an equation means equality under all such
valuations. A proof calculus and an informative axiomatization remain to
be developed.

The preceding results are direct derivations from finite backward
induction. Higher-order games and selection functions already provide
substantial theory for this setting \cite{EscardoOliva2023}.
The calculation above is over trees and sets of winners; it is not a
claim that a monadic product of nondeterministic selection functions
computes all SPE plays. Bolt, Hedges, and Zahn establish important
limitations on precisely that product construction
\cite{BoltHedgesZahn2019}. Establishing an original logical result requires
comparison at these exact interfaces, beyond correctness of the finite
algebra derived here.
